\section{Автономные миссии}
\label{sec:missions}

\begin{infobox}[Цели главы]
\begin{itemize}
    \item Знать жизненный цикл миссии: предстарт $\to$ взлёт $\to$ точки $\to$ посадка.
    \item Отдавать команды автопилоту (\texttt{ap.push}, \texttt{ap.goToLocalPoint}) и дожидаться их выполнения по событиям.
    \item Понимать систему координат сцены и задавать точки в метрах.
    \item Летать по фигурам: прямая, многоугольник, окружность --- точками и таймером.
    \item Строить миссию как конечный автомат с аварийными реакциями и запуском с пульта.
\end{itemize}
\end{infobox}

\subsection{Что такое миссия}

Автономная миссия --- заранее запрограммированная последовательность действий, которую дрон выполняет без участия пилота. Главное, что нужно понять в этой главе: автопилот Пионера --- \textbf{не исполнитель команд по порядку}, а машина событий. Программа \emph{отдаёт команду} («взлетай») и заканчивается. Когда автопилот выполнит команду, он \emph{сообщит об этом событием} (\verb|TAKEOFF_COMPLETE|) --- и только тогда программа отдаёт следующую.

Поэтому программа миссии --- это не список команд, а набор ответов на вопрос «что делать, когда произошло X». Ответы живут в \verb|callback(event)|, и мы уже умеем его писать.

\subsection{Толковый словарь}
\begin{description}
  \item[Автопилот (\texttt{ap})] модуль команд полёта: предстарт, взлёт, посадка, полёт к точке, курс.
  \item[Команда (\texttt{ap.push(Ev.MCE\_*)})] просьба автопилоту сменить режим. Выполняется не мгновенно.
  \item[Событие (\texttt{Ev.*} в \texttt{callback})] сообщение автопилота о том, что команда выполнена или что-то произошло.
  \item[Локальная система координат] оси X, Y на плоскости и Z вверх, начало --- точка старта сцены, единицы --- метры.
  \item[Курс (\texttt{yaw})] угол поворота носа дрона вокруг вертикальной оси, в радианах.
  \item[Счётчик шагов] переменная с номером текущей точки маршрута: по ней программа понимает, какую точку только что достигла.
  \item[Failsafe] аварийная реакция: что делать при ударе, разряде батареи, потере связи.
  \item[RC-каналы (\texttt{Sensors.rc})] значения стиков и тумблеров пульта, от $-1$ до $1$.
\end{description}
